\section{The \acr{GPU} Memory Hierarchy}

\acr{GPU} programming requires understanding the memory hierarchy and synchronization primitives to achieve high performance.

\acr{GPU}s have a complex memory hierarchy with vastly different performance characteristics:

\begin{longtable}{|p{3cm}|p{2.5cm}|p{3cm}|p{4.5cm}|}
\hline
\textbf{Memory Type} & \textbf{Latency} & \textbf{Bandwidth} & \textbf{Use Case} \\
\hline
Registers & 1 cycle & >10 TB/s & Loop counters, temporaries \\
\hline
Shared/\acr{LDS} & 20-30 cycles & 1-2 TB/s & Tile caching, reductions \\
\hline
L1 Cache & 50-100 cycles & 500 GB/s - 1 TB/s & Automatic caching \\
\hline
L2 Cache & 100-200 cycles & 200-500 GB/s & Automatic caching \\
\hline
Global/\acr{VRAM} & 200-400 cycles & 500 GB/s - 2 TB/s & Large datasets \\
\hline
\end{longtable}

The key to \acr{GPU} performance is \textbf{minimizing global memory access} by using registers and shared memory.

